% Fixture for tests/nvim-latency -- self-contained, no repo dependencies.
% Paragraphs are single logical lines of ~600/1000/1400/1900 characters,
% matching the shape of the dissertation chapters this suite guards.
\documentclass[letterpaper,12pt]{article}
\usepackage{french-logic}
\begin{document}
\section{Fixture}

The canonical model for $\CE$ fixes $\fn{w}$ only on the definable propositions, so $\truthset{A}\intersect\truthset{B}$ is again a proof-set while $\proofset{C}$ need not be; by \CMl\ and \CMr\ the frame condition $\cmr$ then quantifies over $\power{W}$ rather than the algebra, and $\M,w\trues A\cnecs B$ holds iff $\eclass{A}\inclin\fn{w}(\truthset{B})$, which is what \cite{chellas1980} calls supplementation. The canonical model for $\CE$ fixes $\fn{w}$ only on the definable propositions, so $\truthset{A}\intersect\truthset{B}$ is again a proof-set while $\proofset{C}$ need not be; by \CMl\ and \CMr\ the frame condition $\cmr$ then quantifies over $\power{W}$ rather than the algebra, and $\M,w\trues A\cnecs B$ holds iff $\eclass{A}\inclin\fn{w}(\truthset{B})$, which is what \cite{chellas1980} calls supplementation.

The canonical model for $\CE$ fixes $\fn{w}$ only on the definable propositions, so $\truthset{A}\intersect\truthset{B}$ is again a proof-set while $\proofset{C}$ need not be; by \CMl\ and \CMr\ the frame condition $\cmr$ then quantifies over $\power{W}$ rather than the algebra, and $\M,w\trues A\cnecs B$ holds iff $\eclass{A}\inclin\fn{w}(\truthset{B})$, which is what \cite{chellas1980} calls supplementation. The canonical model for $\CE$ fixes $\fn{w}$ only on the definable propositions, so $\truthset{A}\intersect\truthset{B}$ is again a proof-set while $\proofset{C}$ need not be; by \CMl\ and \CMr\ the frame condition $\cmr$ then quantifies over $\power{W}$ rather than the algebra, and $\M,w\trues A\cnecs B$ holds iff $\eclass{A}\inclin\fn{w}(\truthset{B})$, which is what \cite{chellas1980} calls supplementation. The canonical model for $\CE$ fixes $\fn{w}$ only on the definable propositions, so $\truthset{A}\intersect\truthset{B}$ is again a proof-set while $\proofset{C}$ need not be; by \CMl\ and \CMr\ the frame condition $\cmr$ then quantifies over $\power{W}$ rather than the algebra, and $\M,w\trues A\cnecs B$ holds iff $\eclass{A}\inclin\fn{w}(\truthset{B})$, which is what \cite{chellas1980} calls supplementation.

The canonical model for $\CE$ fixes $\fn{w}$ only on the definable propositions, so $\truthset{A}\intersect\truthset{B}$ is again a proof-set while $\proofset{C}$ need not be; by \CMl\ and \CMr\ the frame condition $\cmr$ then quantifies over $\power{W}$ rather than the algebra, and $\M,w\trues A\cnecs B$ holds iff $\eclass{A}\inclin\fn{w}(\truthset{B})$, which is what \cite{chellas1980} calls supplementation. The canonical model for $\CE$ fixes $\fn{w}$ only on the definable propositions, so $\truthset{A}\intersect\truthset{B}$ is again a proof-set while $\proofset{C}$ need not be; by \CMl\ and \CMr\ the frame condition $\cmr$ then quantifies over $\power{W}$ rather than the algebra, and $\M,w\trues A\cnecs B$ holds iff $\eclass{A}\inclin\fn{w}(\truthset{B})$, which is what \cite{chellas1980} calls supplementation. The canonical model for $\CE$ fixes $\fn{w}$ only on the definable propositions, so $\truthset{A}\intersect\truthset{B}$ is again a proof-set while $\proofset{C}$ need not be; by \CMl\ and \CMr\ the frame condition $\cmr$ then quantifies over $\power{W}$ rather than the algebra, and $\M,w\trues A\cnecs B$ holds iff $\eclass{A}\inclin\fn{w}(\truthset{B})$, which is what \cite{chellas1980} calls supplementation. The canonical model for $\CE$ fixes $\fn{w}$ only on the definable propositions, so $\truthset{A}\intersect\truthset{B}$ is again a proof-set while $\proofset{C}$ need not be; by \CMl\ and \CMr\ the frame condition $\cmr$ then quantifies over $\power{W}$ rather than the algebra, and $\M,w\trues A\cnecs B$ holds iff $\eclass{A}\inclin\fn{w}(\truthset{B})$, which is what \cite{chellas1980} calls supplementation.

The canonical model for $\CE$ fixes $\fn{w}$ only on the definable propositions, so $\truthset{A}\intersect\truthset{B}$ is again a proof-set while $\proofset{C}$ need not be; by \CMl\ and \CMr\ the frame condition $\cmr$ then quantifies over $\power{W}$ rather than the algebra, and $\M,w\trues A\cnecs B$ holds iff $\eclass{A}\inclin\fn{w}(\truthset{B})$, which is what \cite{chellas1980} calls supplementation. The canonical model for $\CE$ fixes $\fn{w}$ only on the definable propositions, so $\truthset{A}\intersect\truthset{B}$ is again a proof-set while $\proofset{C}$ need not be; by \CMl\ and \CMr\ the frame condition $\cmr$ then quantifies over $\power{W}$ rather than the algebra, and $\M,w\trues A\cnecs B$ holds iff $\eclass{A}\inclin\fn{w}(\truthset{B})$, which is what \cite{chellas1980} calls supplementation. The canonical model for $\CE$ fixes $\fn{w}$ only on the definable propositions, so $\truthset{A}\intersect\truthset{B}$ is again a proof-set while $\proofset{C}$ need not be; by \CMl\ and \CMr\ the frame condition $\cmr$ then quantifies over $\power{W}$ rather than the algebra, and $\M,w\trues A\cnecs B$ holds iff $\eclass{A}\inclin\fn{w}(\truthset{B})$, which is what \cite{chellas1980} calls supplementation. The canonical model for $\CE$ fixes $\fn{w}$ only on the definable propositions, so $\truthset{A}\intersect\truthset{B}$ is again a proof-set while $\proofset{C}$ need not be; by \CMl\ and \CMr\ the frame condition $\cmr$ then quantifies over $\power{W}$ rather than the algebra, and $\M,w\trues A\cnecs B$ holds iff $\eclass{A}\inclin\fn{w}(\truthset{B})$, which is what \cite{chellas1980} calls supplementation. The canonical model for $\CE$ fixes $\fn{w}$ only on the definable propositions, so $\truthset{A}\intersect\truthset{B}$ is again a proof-set while $\proofset{C}$ need not be; by \CMl\ and \CMr\ the frame condition $\cmr$ then quantifies over $\power{W}$ rather than the algebra, and $\M,w\trues A\cnecs B$ holds iff $\eclass{A}\inclin\fn{w}(\truthset{B})$, which is what \cite{chellas1980} calls supplementation.

\end{document}
